\documentclass[12pt]{article}
\usepackage[T1]{fontenc}
\usepackage{jwstproposaltemplate_v6}
\usepackage{graphicx}
\usepackage{amsmath}
\usepackage[hidelinks]{hyperref}
\begin{document}
\scientificjustification

\noindent\textbf{Does the dust around a changing nucleus keep up with it?}
An active galactic nucleus (AGN) can undergo large luminosity changes within years [1,2], sometimes accompanied by a transformation of its continuum and broad-line spectrum [3--5]. Dust responds asymmetrically: when a nucleus brightens, sublimation clears the inner boundary within months [15--18], but after a fade the boundary can only return as fast as dust condenses, a timescale that has never been measured. Light-travel delays add a further lag [6,7]. Changing-look AGN followed through rises, fades and recoveries are the laboratory for both.

We request \textbf{67.0~h of MIRI/MRS spectroscopy of 24 AGN at $z<0.3$: 12 fading and 12 rising sources}. Every target has ZTF/WISE light curves, SPHEREx spectra and multi-epoch optical spectroscopy, including 2026 Palomar/NGPS spectra. Together, their recorded histories provide ensemble coverage of different stages of nuclear activity and dust response. MIRI supplies the essential warm continuum, both silicate features and spatial constraints on host emission, connecting each history to its warm-dust spectrum. The programme addresses two linked questions:

\noindent\textbf{1. Is the dust response hysteretic?} If re-formation is slow, a nucleus that reached its present luminosity by rising carries less warm dust than one that reached it by fading. Twelve rising and twelve fading targets, each anchored to its own 2010 warm-dust measurement, test this; delayed heating of unchanged dust is the null model.

\noindent\textbf{2. Does the dust boundary move, or only its temperature?} Fixed and evolving distributions predict different warm spectra for the same history [6,7,18,19]. The silicate features and the full 5--28~$\mu$m continuum separate the two.

\noindent\textbf{From infrared echoes to dust evolution.}
WISE observations have established infrared echoes of changing-look transitions [8--11], supporting intrinsic changes in nuclear power in those sources. Broadband colours, however, combine temperature, emitting area, extinction and host light [12--14]. The proposed spectra test the warm spectral signature of those histories and constrain which properties of the emitting material produce it.

The equilibrium sublimation radius scales approximately as $L^{1/2}$ [15--17], but the physical dust boundary can adjust over years [18]. Individual nuclei already suggest different responses: the smaller dust radius after Mrk~590 faded supports replenishment [19], whereas the temperature variations in Mrk~817 favour heating and cooling at a nearly fixed boundary [20]. These results motivate a comparison across both directions of change. Recovering sources further test whether the spectrum depends on the path to its present state.

\noindent\textbf{MIRI's diagnostic power.}
SPHEREx covers 0.75--5~$\mu$m [21], constraining hot dust close to the nucleus [20,22]. Longer delays and weaker variability at longer wavelengths [6,23,24] make the warm continuum sensitive to earlier illumination. MRS supplies the missing 4.9--27.9~$\mu$m spectrum [25]: the continuum measures the imprint, the paired silicate features constrain geometry and optical depth [26,27], and spatial profiles and aromatic bands constrain host emission [14,28,29]. NEOWISE and SPHEREx sample only hot dust, so MRS supplies the first warm-dust measurement of these nuclei since AllWISE in 2010 and, with SPHEREx, the full dust temperature ladder within a year.

\clearpage
\begin{figure}[!t]
\centering
\includegraphics[width=\textwidth]{fig1_targets.pdf}
\caption{\textbf{Recorded histories of three sample targets.} Each block: SDSS image (5$''$ bar); ZTF $g$, $r$ nightly medians with alert photometry to 2026 September (open circles); AllWISE and NEOWISE W1, W2; archival SDSS/SDSS-V and 2026 NGPS spectra joined to the 2025--2026 SPHEREx spectrum on one observed-wavelength axis. Light-curve ticks mark the spectral epochs, colour-matched. Top: a flare seen in both ZTF and WISE in 2018--2019, with a much brighter optical continuum in 2019 and a fade since. Middle: a fading case, with WISE declining since 2010 while the optical has brightened modestly. Bottom: a rising case in the light curves and the spectra alike. Dashed: host+disc+blackbody fits to SPHEREx ($T_{\rm col}\approx1100$--1200 K, model-dependent); gold markers: expected lines; no rescaling between epochs.}
\label{fig:history}
\end{figure}

\noindent\textbf{Tracking changes in AGN activity.}
We combine catalogued changing-look AGN [32--35] with large-amplitude ZTF/WISE variables [5,36], emphasizing variability direction, amplitude and timing across spectral classifications [33,37]. A variability manifold groups similar light-curve shapes to identify candidates [38]; multi-epoch spectra establish continuum and broad-line changes. The 12 rising and 12 fading selection histories span sustained changes, plateaus and recoveries. Each target has approximately a decade of W1/W2 monitoring; W1 95th-to-5th-percentile flux ratios span 1.3--2.3 before host subtraction. \textbf{This diversity tests infrared memory across different stages of change.} Recovering sources test whether nuclei at comparable current power retain different warm spectra because they followed different paths. SPHEREx and the 2026 NGPS spectra connect these histories to hot-dust emission and optical spectral state (Figure~\ref{fig:history}).

\clearpage
\noindent\textbf{Measuring dust hysteresis.}
Our primary observable is the nuclear 8--13~$\mu$m flux at the MIRI epoch relative to the 2010 AllWISE measurement of the same source, compared with the flux predicted from each target's own W1 history (Figure~\ref{fig:forecast}). The 2010 point anchors every object, so no equilibrium covering factor is required; the hypotheses differ only through the history: instantaneous response, delayed heating of a fixed distribution [6,7], and an inner boundary that follows $L^{1/2}$ as changed illumination arrives [15--18]. Rising targets should fall below the instantaneous prediction and fading targets above it; their contrast in the sign-weighted residual is the hysteresis measurement, and its dependence on event age bounds the re-formation timescale. Optical decomposition and SPHEREx supply the nuclear fraction of the W1 history [4,20,22]; the silicate profile and 5--20~$\mu$m continuum fix the temperature, and hence the radius and light-travel time, of the dust emitting at 12~$\mu$m [26,28]; aromatic bands and the spatial decomposition give the host contribution to the 2010 aperture [14,29]; optical monitoring near JWST fixes the nuclear state at the MIRI epoch. We will marginalize over event dates, light-curve gaps and the optical-to-heating conversion.

\noindent\textbf{Forecast.} Driving optically thin shells [45] with the measured histories (Figure~\ref{fig:forecast}), the delayed-heating and instantaneous predictions differ by a sign-weighted mean of $0.06\pm0.01$ dex over the 24 targets ($4.5\sigma$), and the evolving-boundary prediction differs from delayed heating by $0.05\pm0.01$ dex ($3.7\sigma$), with per-target errors of 0.03--0.05 dex from the 2010 photometry, the decomposition and the host fraction. Twelve targets separate at least one pair of hypotheses at $3\sigma$ once the host fraction is measured, and the signals grow if the current nuclear trends continue to 2028. Common calibration offsets cancel in the sign-weighted means. Mock MRS spectra spanning the allowed histories and geometries will calibrate false-positive rates and upper limits.

\noindent\textbf{Testing for dust evolution.}
We will use clumpy and disc-plus-wind geometries [40--42] to construct responses to the reconstructed histories [6,7], allowing for composition and grain-size uncertainty [43]. Fixed distributions respond through grain heating and cooling [6]; evolving distributions also change the amount or location of dust [18,19,22]. Both include light-travel delays. The diagnostic $\delta_{\rm warm}=\log_{10}[L^{\rm MRS}_{8\text{--}13\,\mu{\rm m}}/L^{\rm fixed}_{8\text{--}13\,\mu{\rm m}}(t_{\rm MIRI})]$ measures the warm-continuum residual relative to a \emph{delayed fixed-distribution} prediction. The known limitations of static spectral models [28] motivate a joint test: evolving models must explain coherent continuum and silicate residuals across the allowed fixed models while also fitting the earlier data, with host and calibration uncertainties propagated throughout [44].

Figure~\ref{fig:predictions} illustrates both hypotheses for a step change. Where both fit a target, we still measure or bound the delayed-heating term and the change in its warm covering factor since 2010 [30,31].

\noindent\textbf{Mechanism: dimmed or hidden, and an extreme-ultraviolet clock.}
For the fading half, the 9.7~$\mu$m silicate profile stays in emission if the nucleus dimmed and turns to absorption if it became obscured [26,27], and [Ne\,\textsc{v}] 14.3/[Ne\,\textsc{ii}] 12.8~$\mu$m shows whether the ionising continuum fell; the rising half gives the reverse test. The coronal lines [Ne\,\textsc{vi}] 7.7 and [Ne\,\textsc{v}] 14.3~$\mu$m ([Ne\,\textsc{v}] 24.3 and [O\,\textsc{iv}] 25.9~$\mu$m for the twelve targets at $z<0.15$) arise at parsec distances and average the extreme-ultraviolet output over the past few years [50]; their strength relative to [Ne\,\textsc{iii}] 15.6~$\mu$m in rising versus fading nuclei is a second history clock and tests whether changing-look transitions are accretion-state changes [3]. PAH bands and H$_2$ lines, resolved on 0.5--2~kpc, show whether the hosts of the two halves differ [29,49].

\begin{figure}[!p]
\centering
\includegraphics[width=.78\textwidth]{fig2_diagnostics.pdf}
\caption{\textbf{Delayed heating and dust evolution for a step change.} A: equal current power, different histories. B: fixed optically thin dust with light-travel delays [6,45] after a fourfold rise (blue) or decline (red), divided by equilibrium at the current power, 3 and 10 inner-radius light-crossing times later; shading marks common MRS coverage and the 8--13~$\mu$m band. C: the 8--13~$\mu$m offset against event age for twofold and fourfold changes, with the $\pm1\sigma$ precision of a 12-target mean (grey). D: MRS point-spread function at 10~$\mu$m [51] and a 3 kpc exponential host at the sample median and maximum redshift (dashed: channel 2 field [25]). B and C are illustrative, not target predictions.}
\label{fig:predictions}
\includegraphics[width=.83\textwidth]{fig2_forecast.pdf}
\caption{\textbf{What one MIRI epoch measures.} Top left: the W1 (hot dust) history of one target relative to its 2010 AllWISE epoch, and the model nuclear 12~$\mu$m flux under no memory, delayed heating of fixed dust and an evolving inner boundary [6,7,45], anchored to the 2010 AllWISE 12~$\mu$m point, with the adopted MIRI $\pm1\sigma$. Bottom left: its MRS spectrum at the MIRI epoch relative to no memory, scaled at 5~$\mu$m, against the MRS statistical and calibration bands, with the host, ionising-power and geometry diagnostics marked. Right: the three predictions for all 24 targets from their own W1 histories, with the measurement error (dark) and the current host-fraction term (light). Amplitudes are illustrative; sign-weighted means are robust to common calibration offsets.}
\label{fig:forecast}
\end{figure}

\descriptionofobservations

\noindent\textbf{Spectral coverage and sensitivity.}
We request one MIRI/MRS science visit and a matching sky visit per target. All three grating settings, nuclear target acquisition and a four-point dither provide approximately 4.9--27.9~$\mu$m coverage [25,46]. Both silicate peaks fall within this range for all 24 targets. The primary rest-frame 8--13~$\mu$m continuum diagnostic is accessible throughout the sample without channel~4C; the full spectrum constrains the longer-wavelength continuum and 18~$\mu$m profile.

Archival total-source W3 and W4 fluxes span 1.9--48.4 and 6.4--82.6~mJy. We target continuum S/N of 50 per $R=100$ bin near rest 5 and 12~$\mu$m and 30 around the 18~$\mu$m feature, corresponding to 2\% and 3.3\% statistical errors per bin. We use FULL/FASTR1, one integration per dither, and 27/60/27 groups in settings A/B/C for 23 targets. Four dithers give 299.7/666.0/299.7~s on source, with all four channels observed simultaneously. The W3-faint J083826.50+371906.7 uses 80 groups in B (888.0~s). Equal-depth sky exposures preserve extended host emission for decomposition [46]. Science-plus-sky integration totals 17.00~h. \textbf{We request 67.0~h}, including acquisition, visit operations, slews and observatory overheads calculated with APT~2026.5.1: a Medium programme [47].

For J083826.50+371906.7 at $z=0.2111$, ETC/Pandeia~2026.7 predicts S/N of 55/225/61 per $R=100$ bin at rest 5/12/18~$\mu$m. These estimates assume a point source with $F_\nu\propto\lambda^2$, 0.95~mJy in W3 (half the archival total), the adopted sky spectrum, equal-depth sky exposures and extraction radii of $0.3''$/$0.6''$/$0.9''$. They establish statistical sensitivity under the adopted nuclear-flux model; host and calibration uncertainties enter the science inference separately. Acquisition uses F560W for the faintest sources and the neutral-density filter for brighter nuclei, with readouts giving S/N above 50 without saturation over the adopted flux ranges [48].

\noindent\textbf{Connecting the observing epochs.}
The historical baseline gives each MIRI spectrum its temporal context, and the 2010 AllWISE W3/W4 photometry anchors each nucleus before its recorded change. WISE and SPHEREx constrain the dust response at their observed epochs; optical spectroscopy tracks the nuclear spectrum. Optical monitoring near JWST anchors the nuclear continuum at the MIRI epoch. One MIRI visit per target connects recorded activity to the warm-dust response across 24 histories spanning different directions, amplitudes and event ages.

\noindent\textbf{Separating nuclear dust from the host.}
\textbf{Each MRS cube measures the nucleus and surrounding host together.} Wavelength-dependent point-spread-function fitting separates unresolved and extended emission [28]; surrounding spaxels constrain the host spectrum directly. Joint modelling distinguishes continuum, silicate, PAH and H$_2$ contributions [13,49]. The 6.2, 7.7 and 11.3~$\mu$m PAH bands constrain star formation with allowance for AGN processing [29]; [Ne\,\textsc{v}] and [Ne\,\textsc{vi}] trace highly ionised gas [50] at S/N above 10 per resolution element. Alternative extractions and host templates [14] define conservative nuclear-flux bounds.

We will compare the 2010 W3/W4 photometry with redshifted nuclear-plus-host models integrated through the WISE bandpasses, including colour corrections and emission outside MRS coverage. Archival imaging constrains host light beyond the MRS field and within the SPHEREx aperture. Both science tests propagate nuclear-flux bounds and wavelength-dependent calibration covariance [44].

\clearpage
\supplementalinformation
\noindent\textbf{Special requirements.} Each science/sky pair is linked as a consecutive, noninterruptible sequence to obtain a contemporaneous background [46]. Sky fields are selected using WISE and 2MASS catalogues and images, with at least $25''$ clearance from catalogued sources. Twenty-three fields lie within 1--2~arcmin of their nuclei; the crowded field around J20145928+2523010 uses a $135''$ offset. Background targets are marked extended and use four-point extended-source dithers, with no target acquisition. No fixed-date constraints, coordinated parallels or joint-facility time are requested. APT's Visit Planner returns scheduling windows for all 48 visits.

\noindent\textbf{Duplications.} A coordinate-based MAST search on 2026 September 26, including archived and planned JWST observations, found no overlapping observations within $120''$ of any of the 24 science positions. No potential MIRI/MRS duplications were identified.

\references
\noindent [1] \href{https://doi.org/10.3847/1538-4357/aaa9b6}{Rumbaugh, N., et al. 2018, ApJ, 854, 160.}\par
\noindent [2] \href{https://doi.org/10.1093/mnras/stz3244}{Graham, M. J., Ross, N. P., Stern, D., et al. 2020, MNRAS, 491, 4925.}\par
\noindent [3] \href{https://www.nature.com/articles/s41550-023-02108-4}{Ricci, C., \& Trakhtenbrot, B. 2023, Nature Astronomy, 7, 1282.}\par
\noindent [4] \href{https://doi.org/10.1088/0004-637X/800/2/144}{LaMassa, S. M., et al. 2015, ApJ, 800, 144.}\par
\noindent [5] \href{https://doi.org/10.1093/mnras/stv2997}{MacLeod, C. L., et al. 2016, MNRAS, 457, 389.}\par
\noindent [6] \href{https://doi.org/10.1051/0004-6361/201117750}{H\"onig, S. F., \& Kishimoto, M. 2011, A\&A, 534, A121.}\par
\noindent [7] \href{https://doi.org/10.3847/1538-4357/ab6aa1}{Almeyda, T., et al. 2020, ApJ, 891, 26.}\par
\noindent [8] \href{https://authors.library.caltech.edu/records/qjtjh-0cw44}{Sheng, Z., et al. 2017, ApJL, 846, L7.}\par
\noindent [9] \href{https://authors.library.caltech.edu/records/azy2b-r2w44}{Stern, D., et al. 2018, ApJ, 864, 27.}\par
\noindent [10] \href{https://arxiv.org/abs/1711.08122}{Yang, Q., et al. 2018, ApJ, 862, 109.}\par
\noindent [11] \href{https://doi.org/10.1093/mnras/sty2002}{Ross, N. P., Ford, K. E. S., Graham, M., et al. 2018, MNRAS, 480, 4468.}\par
\noindent [12] \href{https://doi.org/10.3390/universe8060304}{Lyu, J., \& Rieke, G. 2022, Universe, 8, 304.}\par
\noindent [13] \href{https://doi.org/10.1093/mnras/stae612}{Donnan, F. R., et al. 2024, MNRAS, 529, 1386.}\par
\noindent [14] \href{https://doi.org/10.1088/0004-637X/803/2/109}{Hern\'an-Caballero, A., et al. 2015, ApJ, 803, 109.}\par
\noindent [15] \href{https://articles.adsabs.harvard.edu/pdf/1987ApJ...320..537B}{Barvainis, R. 1987, ApJ, 320, 537.}\par
\noindent [16] \href{https://doi.org/10.1051/0004-6361:20077911}{Kishimoto, M., H\"onig, S. F., Beckert, T., \& Weigelt, G. 2007, A\&A, 476, 713.}\par
\noindent [17] \href{https://arxiv.org/abs/1406.2078}{Koshida, S., et al. 2014, ApJ, 788, 159.}\par
\noindent [18] \href{https://arxiv.org/abs/1308.6517}{Kishimoto, M., et al. 2013, ApJL, 775, L36.}\par
\noindent [19] \href{https://academic.oup.com/mnras/article/491/4/4615/5652205}{Kokubo, M., \& Minezaki, T. 2020, MNRAS, 491, 4615.}\par
\noindent [20] \href{https://arxiv.org/abs/2510.21395}{Landt, H., et al. 2026, ApJ, 997, 22.}\par
\noindent [21] IPAC/IRSA, \href{https://irsa.ipac.caltech.edu/onlinehelp/spherex/spherex/overview.html}{``SPHEREx Data Explorer: Overview''}.\par
\noindent [22] \href{https://arxiv.org/abs/1908.01627}{Landt, H., et al. 2019, MNRAS, 489, 1572.}\par
\noindent [23] \href{https://doi.org/10.3847/1538-4357/abee14}{Lyu, J., \& Rieke, G. H. 2021, ApJ, 912, 126.}\par
\noindent [24] \href{https://doi.org/10.3847/1538-4357/ab481d}{Lyu, J., Rieke, G. H., \& Smith, P. S. 2019, ApJ, 886, 33.}\par
\noindent [25] STScI, JWST User Documentation, \href{https://jwst-docs.stsci.edu/jwst-mid-infrared-instrument/miri-observing-modes/miri-medium-resolution-spectroscopy}{``MIRI Medium Resolution Spectroscopy''}.\par
\noindent [26] \href{https://doi.org/10.1088/0004-637X/697/1/182}{Thompson, G. D., Levenson, N. A., Uddin, S. A., \& Sirocky, M. M. 2009, ApJ, 697, 182.}\par
\noindent [27] \href{https://doi.org/10.1086/586727}{Sirocky, M. M., et al. 2008, ApJ, 678, 729.}\par
\noindent [28] \href{https://arxiv.org/abs/2504.01103}{Gonz\'alez-Mart\'in, O., et al. 2025, MNRAS, 539, 2158.}\par
\noindent [29] \href{https://doi.org/10.1051/0004-6361/202450086}{Garc\'ia-Bernete, I., et al. 2024, A\&A, 691, A162.}\par
\noindent [30] \href{https://doi.org/10.1093/mnras/stw444}{Stalevski, M., et al. 2016, MNRAS, 458, 2288.}\par
\noindent [31] \href{https://doi.org/10.1093/mnras/stv299}{Mateos, S., et al. 2015, MNRAS, 449, 1422.}\par
\noindent [32] \href{https://arxiv.org/abs/2206.10056}{Panda, S., \& \'Sniegowska, M. 2024, ApJS, 272, 13.}\par
\noindent [33] \href{https://arxiv.org/abs/2401.01933}{Zeltyn, G., et al. 2024, ApJ, 966, 85.}\par
\noindent [34] \href{https://arxiv.org/abs/2408.00402}{Guo, W.-J., et al. 2025, ApJS, 278, 28.}\par
\noindent [35] \href{https://arxiv.org/abs/2408.07335}{Dong, Q., et al. 2025, ApJ, 986, 160.}\par
\noindent [36] \href{https://doi.org/10.3847/1538-4357/ab5af9}{Sheng, Z., et al. 2020, ApJ, 889, 46.}\par
\noindent [37] \href{https://doi.org/10.3847/1538-4357/ae1fdc}{Wang, H., et al. 2026, ApJ, 997, 100.}\par
\noindent [38] \href{https://doi.org/10.3847/1538-4357/ae38b8}{Hemmati, S., et al. 2026, ApJ, 998, 130.}\par
\noindent [39] \href{https://doi.org/10.3847/1538-4357/ab4f7b}{Minezaki, T., et al. 2019, ApJ, 886, 150.}\par
\noindent [40] \href{https://arxiv.org/abs/0806.0512}{Nenkova, M., et al. 2008, ApJ, 685, 160.}\par
\noindent [41] \href{https://arxiv.org/abs/1306.4312}{H\"onig, S. F., Kishimoto, M., Tristram, K. R. W., et al. 2013, ApJ, 771, 87.}\par
\noindent [42] \href{https://doi.org/10.3847/2041-8213/aa6838}{H\"onig, S. F., \& Kishimoto, M. 2017, ApJL, 838, L20.}\par
\noindent [43] \href{https://doi.org/10.1093/mnras/staf1494}{Reyes-Amador, O. U., et al. 2025, MNRAS, 543, 813.}\par
\noindent [44] STScI, JWST User Documentation, \href{https://jwst-docs.stsci.edu/jwst-calibration-status/miri-calibration-status/miri-mrs-calibration-status}{``MIRI MRS Calibration Status''}.\par
\noindent [45] \href{https://doi.org/10.1086/172149}{Laor, A., \& Draine, B. T. 1993, ApJ, 402, 441}; optical-property tables supplied by B. T. Draine.\par
\noindent [46] STScI, JWST User Documentation, \href{https://jwst-docs.stsci.edu/jwst-mid-infrared-instrument/miri-observing-strategies/miri-mrs-recommended-strategies}{``MIRI MRS Recommended Strategies''} and \href{https://jwst-docs.stsci.edu/jwst-mid-infrared-instrument/miri-operations/miri-dithering/miri-mrs-dedicated-sky-observations}{``MIRI MRS Dedicated Sky Observations''}.\par
\noindent [47] STScI, \href{https://jwst-docs.stsci.edu/jwst-opportunities-and-policies/jwst-call-for-proposals-for-cycle-6}{``JWST Call for Proposals for Cycle 6''} and \href{https://jwst-docs.stsci.edu/jwst-opportunities-and-policies/jwst-call-for-proposals-for-cycle-6/jwst-proposal-types-and-categories/jwst-observations-needing-special-requirements}{``JWST Observations Needing Special Requirements''}.\par
\noindent [48] STScI, JWST User Documentation, \href{https://jwst-docs.stsci.edu/jwst-mid-infrared-instrument/miri-operations/miri-target-acquisition/miri-mrs-target-acquisition}{``MIRI MRS Target Acquisition''} and \href{https://jwst-docs.stsci.edu/jwst-exposure-time-calculator-overview/jwst-etc-calculations-page-overview/jwst-etc-target-acquisition/jwst-etc-miri-target-acquisition}{``JWST ETC MIRI Target Acquisition''}.\par
\noindent [49] \href{https://academic.oup.com/mnras/article/519/3/3691/6948335}{Donnan, F. R., et al. 2023, MNRAS, 519, 3691.}\par
\noindent [50] \href{https://doi.org/10.1051/0004-6361/202452437}{Hermosa Mu\~noz, L., et al. 2025, A\&A, 693, A321.}\par
\noindent [51] \href{https://doi.org/10.3847/1538-3881/acdddc}{Law, D. R., et al. 2023, AJ, 166, 45.}\par

\noindent\textbf{Generative AI disclosure.} OpenAI Codex (GPT-6; used 2026 September 26--28) assisted with editing and literature searches.\par
\templateversion
\end{document}
